\chapter{Schema Integration and Analytical Storage}

After schema matching and entity resolution, the project materializes a unified
restaurant dataset. This stage is the point at which the three platform-specific
clean collections become a single integrated database suitable for quality
assessment and analytical querying.

\section{Construction of the Integrated Dataset}

The integration process is implemented as a MongoDB-to-MongoDB transformation. It
reads the three clean collections and the entity-resolution candidates:

\begin{itemize}
    \item \texttt{restaurants\_clean\_google};
    \item \texttt{restaurants\_clean\_tripadvisor};
    \item \texttt{restaurants\_clean\_thefork};
    \item \texttt{entity\_resolution\_candidates}.
\end{itemize}

The output consists of two collections:

\begin{itemize}
    \item \texttt{entity\_resolution\_links}, containing the selected one-to-one
    links between Google anchors and matched source records;
    \item \texttt{restaurants\_integrated}, containing the final integrated
    restaurant entities.
\end{itemize}

The effective candidate label is computed as:

\[
\texttt{effective\_label} =
\begin{cases}
\texttt{llm\_label}, & \text{if an LLM decision is available}\\
\texttt{label}, & \text{otherwise}
\end{cases}
\]

Only candidates whose effective label is \texttt{MATCH} are eligible for link
selection. Tripadvisor and TheFork links are selected independently, so an
integrated restaurant may contain only Google evidence, Google plus Tripadvisor,
Google plus TheFork, or all three platforms.

\section{Final Integrated Schema}

The integrated collection is Google-seeded: each document represents one dining
and operational Google restaurant. The document identifier is deterministic:

\[
\texttt{integrated\_restaurant\_id} = \texttt{"google:"} + \texttt{place\_id}
\]

Google is used as the canonical identity and geographic backbone because it has
the highest spatial coverage and complete coordinates. Non-Google sources are
attached as evidence blocks only when selected entity-resolution links exist.

\begin{table}[H]
\centering
\footnotesize
\begin{tabularx}{\textwidth}{p{0.28\textwidth}X}
\toprule
\textbf{Field group} & \textbf{Main attributes} \\
\midrule
Canonical identity and geography &
\texttt{integrated\_restaurant\_id}, \texttt{canonical\_name},
\texttt{canonical\_address}, \texttt{canonical\_city},
\texttt{latitude}, \texttt{longitude}. \\
Platform membership &
\texttt{has\_google}, \texttt{has\_tripadvisor}, \texttt{has\_thefork},
\texttt{platform\_count}. \\
Comparable ratings &
\texttt{google\_rating\_5}, \texttt{tripadvisor\_rating\_5},
\texttt{thefork\_rating\_5}, \texttt{rating\_avg\_5},
\texttt{rating\_range\_5}. \\
Review counts &
\texttt{google\_review\_count}, \texttt{tripadvisor\_review\_count}, \texttt{thefork\_review\_count}. \\
Contacts and price &
\texttt{website}, \texttt{phones}, \texttt{price\_level},
\texttt{price\_evidence}. \\
Audit and provenance &
\texttt{integration\_flags}, \texttt{sources.google},
\texttt{sources.tripadvisor}, \texttt{sources.thefork}. \\
\bottomrule
\end{tabularx}
\caption{Structure of the final \texttt{restaurants\_integrated} collection.}
\label{tab:integrated-schema-overview}
\end{table}

The schema deliberately preserves source-specific ratings rather than replacing
them with a single resolved value. This is a central design decision: the research
questions require measuring disagreement across platforms, therefore the
differences between Google, Tripadvisor, and TheFork must remain visible.

\section{Conflict Handling and Provenance}

Conflicts are resolved according to the role of each attribute. Canonical identity
and coordinates use a metadata-based strategy: Google is trusted as the anchor
source. Comparative values, such as ratings and review counts, are not resolved
away; they are kept side by side. Derived fields such as \texttt{rating\_avg\_5}
and \texttt{rating\_range\_5} summarize the comparison without hiding the source
values.

Contacts and price are handled differently. Phone values are preserved as a
deduplicated set because a restaurant may legitimately have multiple numbers.
Website conflicts keep the raw evidence and prefer Google only when conflicting
values cannot be reconciled. Price evidence from Google, Tripadvisor, and TheFork
is mapped to a shared representation, while the original source-specific values
remain available under the source blocks.

This design maintains reproducibility: every integrated top-level value can be
traced back to the original platform evidence stored under
\texttt{sources.google}, \texttt{sources.tripadvisor}, and
\texttt{sources.thefork}.

\section{Post-Integration Quality Assessment}

The quality of the integration was evaluated against hand-labelled gold data. The
assessment measures not only the entity-resolution classifier, but also the
survival of one-to-one selected links and the spatial enrichment of Tripadvisor
records.

\begin{table}[H]
\centering
\footnotesize
\begin{tabular}{p{0.55\textwidth}r}
\toprule
\textbf{Metric} & \textbf{Value} \\
\midrule
Out-of-sample evaluation rows & 288 \\
In-calibration labelled rows & 600 \\
MATCH precision & 96.8\% \\
Strict MATCH recall & 71.7\% \\
MATCH-or-UNCERTAIN kept recall & 100.0\% \\
Uncertain rate & 33.0\% \\
Selected links from human-confirmed matches & 88 / 127 \\
Link survival rate & 69.3\% \\
Integration survival rate & 69.3\% \\
Median Tripadvisor geocoding error & 10.1 m \\
Tripadvisor coordinates within 50 m & 91.3\% \\
Tripadvisor coordinate coverage & 84.0\% \\
\bottomrule
\end{tabular}
\caption{Post-integration assessment of entity resolution, link selection, and Tripadvisor geocoding.}
\label{tab:post-integration-quality}
\end{table}

The results show a conservative integration policy. Precision is high, meaning
that automatic matches are rarely wrong. Strict recall is lower because many
true matches are intentionally left as \texttt{UNCERTAIN} rather than forced into
unsafe automatic links. This behaviour is desirable in a data integration
setting: false positives would contaminate the final dataset more severely than
unresolved candidates.

The survival rate quantifies the effect of one-to-one selection. Some
human-confirmed matches are not present in the final integrated source blocks
because conservative link selection removes ambiguous alternatives. The geocoding
assessment confirms that Tripadvisor spatial enrichment is reliable enough for
blocking and spatial analysis: the median error is 10.1 m and all evaluated gold
rows are within 100 m.

\section{ClickHouse Analytical Layer}

MongoDB is well suited for raw and integrated document storage, but the final
research questions require repeated aggregations, pairwise rating comparisons,
ranking, grouping by geographic area, and correlation-style summaries. For this
reason, the project adds ClickHouse as a column-oriented analytical DBMS.

The ClickHouse loader reads the cleaned and integrated MongoDB collections and
creates four flat analytical tables:

\begin{itemize}
    \item \texttt{restaurants\_integrated};
    \item \texttt{restaurants\_clean\_google};
    \item \texttt{restaurants\_clean\_tripadvisor};
    \item \texttt{restaurants\_clean\_thefork}.
\end{itemize}

The load process is reproducible and idempotent: each run creates the destination
tables if needed, truncates them, and reloads the current MongoDB state. The
integrated table keeps the analytical columns needed for the research questions,
including platform ratings, review counts, platform membership flags, normalized
price tier, cuisine labels, photo counts, and geographic coordinates. Heavy nested
evidence remains in MongoDB, while ClickHouse stores a flat projection optimized
for SQL analysis.

The design therefore separates responsibilities:

\begin{itemize}
    \item MongoDB stores raw, clean, integrated, and auditable document data;
    \item ClickHouse stores flat analytical tables for fast, reproducible queries;
    \item the notebook reads from ClickHouse without modifying either database.
\end{itemize}

This two-layer architecture satisfies the project requirement of modelling and
storing integrated data in DBMSs while also providing an efficient surface for the
final research-question analysis.
